% !TeX root = ../ADCS & GNC Documentation.tex
\chapter{Simulation vs.\ Flight Software Implementation}

\section{Purpose}

The flight software (FSW) and simulation are designed to be algorithmically identical, ensuring that all control, estimation, and guidance behavior observed in simulation directly translates to flight.

Differences between the two implementations arise solely from:
\begin{itemize}
	\item hardware interfaces
	\item timing and execution constraints
	\item real sensor and actuator behavior
\end{itemize}

No intentional changes are made to control laws, state definitions, or estimation algorithms.

\section{Algorithmic Consistency}

All core ADCS components are shared conceptually between simulation and flight:
\begin{itemize}
	\item Attitude representation (quaternions, scalar-last convention)
	\item Error computation (quaternion error definition and hemisphere handling)
	\item Control law (LQR / PD structure using quaternion vector error and body rates)
	\item State definitions and units
\end{itemize}

This ensures:
\begin{itemize}
	\item identical closed-loop behavior in the absence of disturbances
	\item direct validation of controller performance in simulation
\end{itemize}

\section{Key Implementation Differences}

\subsection{Sensor Inputs}

\textbf{Simulation:}
\begin{itemize}
	\item Perfect or modeled sensor data (e.g., noiseless or parametrically noisy IMU, star tracker)
	\item Perfectly predictable intervals between data packets
\end{itemize}

\textbf{Flight:}
\begin{itemize}
	\item Real sensor measurements with:
	\begin{itemize}
		\item noise
		\item bias
		\item misalignment
		\item dropouts and latency
	\end{itemize}
\end{itemize}

\textbf{Implication:}  
The estimation pipeline (filters, normalization, validation) becomes critical in flight but is simplified in simulation.

\subsection{Actuator Behavior}

\textbf{Simulation:}
\begin{itemize}
	\item Ideal torque application or simplified actuator models
\end{itemize}

\textbf{Flight:}
\begin{itemize}
	\item Real actuators with:
	\begin{itemize}
		\item saturation limits
		\item rate limits
		\item deadbands
		\item delays
	\end{itemize}
\end{itemize}

\textbf{Implication:}  
Commands generated by the controller may not be realized exactly, requiring robustness to actuator constraints.

\subsection{Timing and Execution}

\textbf{Simulation:}
\begin{itemize}
	\item Deterministic, fixed-step integration
	\item Synchronous updates between modules
\end{itemize}

\textbf{Flight:}
\begin{itemize}
	\item Discrete-time execution with:
	\begin{itemize}
		\item task scheduling
		\item variable latency
		\item asynchronous sensor updates
	\end{itemize}
\end{itemize}

\textbf{Implication:}  
The controller operates under sampled-data conditions and must tolerate timing jitter.

\subsection{Numerical and Software Environment}

\textbf{Simulation:}
\begin{itemize}
	\item High-performance computing environment (desktop CPU/GPU)
	\item Double precision typically available
\end{itemize}

\textbf{Flight:}
\begin{itemize}
	\item Embedded processor constraints:
	\begin{itemize}
		\item limited compute resources
		\item possible reduced precision
		\item strict real-time requirements
	\end{itemize}
\end{itemize}

\textbf{Implication:}  
Algorithms must remain computationally efficient and numerically stable.

\section{Practical Considerations for Equivalence}

To maintain consistency between simulation and flight:
\begin{itemize}
	\item Identical mathematical formulations are used for all algorithms
	\item Frame definitions and conventions are strictly preserved
	\item Quaternion normalization and hemisphere policies are enforced consistently
	\item Interfaces are designed so that simulated messages match flight message structures
\end{itemize}

Where possible, simulation components are structured to mirror flight software interfaces, enabling direct reuse or minimal translation. However, due to differences in software architecture, messaging systems, and operating environments, some implementations differ structurally between simulation and flight software. These differences do not affect the underlying algorithms or their resulting behavior, and are limited to interface and execution considerations.

\section{Summary}

The simulation environment serves as a high-fidelity representation of the flight system, with the same underlying ADCS algorithms.

Differences between simulation and flight are not algorithmic, but stem from:
\begin{itemize}
	\item real-world sensor and actuator behavior
	\item timing and execution constraints
	\item embedded system limitations
\end{itemize}

As a result, simulation performance is expected to closely match on-orbit behavior, provided these non-ideal effects are properly accounted for.